\documentclass[10pt,a4paper]{amsart}
\usepackage{amsmath,amssymb,amsthm,graphicx}
\usepackage{hyperref}
\numberwithin{equation}{section}

%\usepackage{subfigure}
\usepackage{upgreek}
\usepackage{mathtools}
\usepackage{tikz}
\usepackage[all]{xy}
\usetikzlibrary{matrix,calc,arrows.meta,positioning,patterns}


\newcommand{\bC}{{\mathbb C}}
\newcommand{\bP}{{\mathbb P}}
\newcommand{\bQ}{{\mathbb Q}}
\newcommand{\bR}{{\mathbb R}}
\newcommand{\bZ}{{\mathbb Z}}
\newcommand{\bN}{{\mathbb N}}
\newcommand{\bS}{{\mathbb S}}
\newcommand{\bF}{{\mathbb F}}
\newcommand{\bT}{{\mathbb T}}
\newcommand{\bK}{{\mathbb K}}
\newcommand{\bfk}{{\mathbf k}}
\newcommand{\bfq}{{\mathbf q}}


\newcommand{\cF}{\mathcal F}
\newcommand{\cG}{\mathcal G}
\newcommand{\cA}{\mathcal A}
\newcommand{\cB}{\mathcal B}
\newcommand{\cC}{\mathcal C}
\newcommand{\cD}{\mathcal D}
\newcommand{\cE}{\mathcal E}
\newcommand{\cK}{\mathcal{K}}
\newcommand{\cY}{\mathcal Y}
\newcommand{\cS}{\mathcal S}
\newcommand{\cN}{\mathcal N}
\newcommand{\cM}{\mathcal M}
\newcommand{\cO}{\mathcal O}
\newcommand{\cP}{\mathcal P}
\newcommand{\cU}{\mathcal U}
\newcommand{\cL}{\mathcal L}
\newcommand{\cR}{\mathcal R}
\newcommand{\cT}{\mathcal T}
\newcommand{\cH}{\mathcal H}
\newcommand{\oD}{\overline{D}}
\newcommand{\scrS}{\EuScript S}
\newcommand{\scrC}{\EuScript C}
\newcommand{\scrD}{\EuScript D}
\newcommand{\scrE}{\EuScript E}
\newcommand{\scrF}{\EuScript F}
\newcommand{\scrG}{\EuScript G}
\newcommand{\scrL}{\EuScript L}
\newcommand{\scrM}{\EuScript M}
\newcommand{\scrO}{\EuScript O}
\newcommand{\scrP}{\EuScript P}
\newcommand{\scrU}{\EuScript U}
\newcommand\isomto{\stackrel{\sim}{\smash{\longrightarrow}\rule{0pt}{0.4ex}}}
\newcommand{\ev}{\operatorname{ev}}
\DeclareMathOperator{\val}{val}
\DeclareMathOperator{\skel}{Skel}
\DeclareMathOperator{\id}{id}
\DeclareMathOperator{\im}{im}
\DeclareMathOperator{\hoSC}{hoSC}
\DeclareMathOperator{\whoSC}{\widehat{hoSC}}
\DeclareMathOperator{\SC}{SC}
\DeclareMathOperator{\wSC}{\widehat{SC}}
\DeclareMathOperator{\wtel}{\widehat{tel}}
\DeclareMathOperator{\Sp}{span}
\DeclareMathOperator{\tel}{tel}
\DeclareMathOperator{\inj}{inj}
\DeclareMathOperator{\conv}{conv}
\DeclareMathOperator{\Sec}{sec}
\DeclareMathOperator{\Area}{Area}
\DeclareMathOperator{\Vol}{Vol}
\DeclareMathOperator{\diam}{diam}
\DeclareMathOperator{\Tot}{Tot}
\newtheorem{Theorem}{Theorem}[section]
\newtheorem*{Theorem*}{Theorem}
\newtheorem{Corollary}[Theorem]{Corollary}
\newtheorem{Lemma}[Theorem]{Lemma}
\newtheorem{tm}[Theorem]{Theorem}
\newtheorem{lm}[Theorem]{Lemma}
\newtheorem{prp}[Theorem]{Proposition}
\newtheorem{Proposition}[Theorem]{Proposition}
\newtheorem{cy}[Theorem]{Corollary}

\theoremstyle{definition}
\newtheorem{Definition}[Theorem]{Definition}
\newtheorem{Note}[Theorem]{Note}
\newtheorem{Example}[Theorem]{Example}
\newtheorem{Remark}[Theorem]{Remark}
\newtheorem{df}[Theorem]{Definition}
\newtheorem{ques}[Theorem]{Question}
\newtheorem{ass}[Theorem]{Assumption}
\newtheorem{obs}[Theorem]{Observation}
\newtheorem{ex}[Theorem]{Example}
\newtheorem{rem}[Theorem]{Remark}
\usepackage{cleveref}
\crefname{originalreference}{original reference}{original references}
\Crefname{originalreference}{Original reference}{Original references}
\makeatletter
\newlabel{prop:starshape-hypersurface}{{1.22}{source}{}{sourcepointer-1}{}}
\newlabel{prop:starshape-hypersurface@cref}{{[originalreference][0][]1.22}{[1][0][]source}{}{}{}}
\makeatother
\crefname{originalreference}{original reference}{original references}
\Crefname{originalreference}{Original reference}{Original references}
\makeatletter
\newlabel{rem:omega-tame}{{3.10}{source}{}{sourcepointer-2}{}}
\newlabel{rem:omega-tame@cref}{{[originalreference][0][]3.10}{[1][0][]source}{}{}{}}
\makeatother
\crefname{originalreference}{original reference}{original references}
\Crefname{originalreference}{Original reference}{Original references}
\makeatletter
\newlabel{tm:normalized-completeness}{{1.16}{source}{}{sourcepointer-3}{}}
\newlabel{tm:normalized-completeness@cref}{{[originalreference][0][]1.16}{[1][0][]source}{}{}{}}
\makeatother
\title{Symplectic hypersurface excision obstructs geometric boundedness}
\author{Yoel Groman}
\date{doi:10.3842/SIGMA.2026.066; CC BY 4.0}
\begin{document}
\maketitle
\noindent\textit{Allowed standard prerequisites (editorial).}\begin{itemize}
\item Symplectic cutting, holomorphic-curve compactness and standard toric/Fukaya background as used by author. Original Sikorav, STVFluxShape and FOOO/ChoOh bibliography retained; actual source-local arguments included.
\end{itemize}\medskip
\ifdefined\hypertarget\hypertarget{sourcepointer-1}{}\fi
\noindent\textit{Source reference.} Original reference 1.22: algebra-geometry/sources/sigma-2026-066/native/sigma26-066.tex, source lines 381--381. Later application or generalization mentioned in the author narrative; not an input to the retained target proof.\par\medskip
\ifdefined\hypertarget\hypertarget{sourcepointer-2}{}\fi
\noindent\textit{Source reference.} Original reference 3.10: algebra-geometry/sources/sigma-2026-066/native/sigma26-066.tex, source lines 998--998. Later application or generalization mentioned in the author narrative; not an input to the retained target proof.\par\medskip
\ifdefined\hypertarget\hypertarget{sourcepointer-3}{}\fi
\noindent\textit{Source reference.} Original reference 1.16: algebra-geometry/sources/sigma-2026-066/native/sigma26-066.tex, source lines 306--306. Later application or generalization mentioned in the author narrative; not an input to the retained target proof.\par\medskip
\subsection{The results}
As our first proxy for completeness, we use the well-established notion of \emph{geometric boundedness} (see, e.g., \cite[Section~4.2]{Sikorav}): an~$\omega$-compatible almost complex structure $J$ on $M$ is called \emph{geometrically bounded} if the associated Riemannian metric $g_J(\cdot,\cdot)=\omega(\cdot,J\cdot)$ is complete and has uniformly bounded sectional curvature and uniformly positive injectivity radius.\footnote{To treat $\omega$-tame almost complex structures, one must also impose a~uniform bound on the norm of $\omega$ with respect to $g_J$ (i.e., $|\omega|_{g_J}\le C$).}

We begin with an existence question.
\begin{ques}\label{ques:1}
Given $\omega$, does there exist a~geometrically bounded $\omega$-compatible almost complex structure $J$ on $M$?
\end{ques}
It is customary to call $(M,\omega)$ itself geometrically bounded if the answer to Question~\ref{ques:1} is affirmative. We consider this question when we start with a~geometrically bounded $M$ and excise a~subset $N\subset M$. Question~\ref{ques:1} can be understood as the question: \emph{can the distance to $N$ of any compact subset of $M\setminus N$ be made arbitrarily large by adjusting the almost complex structure?}




\begin{tm}\label{tm:1}
Let $M$ be a~geometrically bounded symplectic manifold with finitely generated spherical homology, and suppose $\omega$ is rational.
Let $N\subset M$ be a~properly embedded symplectic hypersurface.
Then $M\setminus N$ is not geometrically bounded.
\end{tm}

\section{Proofs}\label{sec:proofs}

\subsection[Energy confinement for J-holomorphic curves]{Energy confinement for $\boldsymbol{J}$-holomorphic curves}
Throughout, if $u\colon (\Sigma,\partial\Sigma)\to (M,\omega)$ is a~$J$-holomorphic curve with (possibly empty) boundary, its energy is
$
E(u):=\int_{\Sigma} u^{*}\omega \ge 0$.


The following confinement lemmas are stated for $\omega$-compatible $J$. At the end of this subsection we note that both hold for $\omega$-tame $J$ under the uniform tameness condition from the footnote to the definition of geometric boundedness on page \pageref{tm:1} (in Section 1.1).


\begin{lm}[Sikorav's confinement estimate]\label{lem:sikorav}
 Let $(M,\omega)$ be a~symplectic manifold and let $J$ be an $\omega$-compatible almost complex structure such that the associated metric $g_J(\cdot,\cdot):=\omega(\cdot,J\cdot)$ is complete and has uniformly bounded sectional curvature and positive injectivity radius $($i.e., $J$ is geometrically bounded$)$.
 Then there exist a~constant $r_0>0$ depending on the geometry and a~universal constant $c>0$ such that for every connected $J$-holomorphic curve $u$ one has
 \[
 u(\Sigma)\subset B_{r_0\,\lceil E(u)/(c r_0^2)\rceil}(u(\partial\Sigma)),
 \]
 where $B_{r}(\cdot)$ denotes the closed $g_J$-ball of radius $r$ about a~subset and $\lceil\cdot\rceil$ is the ceiling function.
 \end{lm}
 \begin{proof}
 See \cite[Proposition 4.4.1]{Sikorav}.
 \end{proof}
 This confinement estimate is concerned with the large scale. Namely, it contains no information about curves with energy much smaller than $cr_0^2$ other than that the diameter is $\leq r_0$. In some cases, e.g., in the proof of Proposition~\ref{prop:starshape-hypersurface}, we need an estimate for the small scale. This is taken care of by the following lemma.
 \begin{lm}[improved confinement]\label{lem:sikorav-improved}
 In the setting of Lemma~{\rm\ref{lem:sikorav}}, with the same $r_0$ and $c$, and letting $C := 1/c$, one has
 \[
 u(\Sigma)\subset B_{\rho}(u(\partial\Sigma)),\qquad \rho:=\max\bigl(3CE(u)/r_0,\sqrt{CE(u)}\bigr).
 \]
 \end{lm}
 \begin{proof}
 The argument relies on the same monotonicity inequality used for Sikorav's lemma: if a~$J$-holomorphic curve meets the center of a~ball of radius $r\le r_0$ and the boundary of the $J$-holomorphic curve does not meet the interior of the ball, then the energy in that ball is at least $c r^2$.

 If $E(u)/\bigl(c r_0^2\bigr)\le 1$, we apply the monotonicity inequality directly: if the image contained a~point at the center of a~ball of radius greater than \smash{$\sqrt{E(u)/c}$} whose interior did not meet $u(\partial\Sigma)$, then monotonicity would force $E(u)\ge c r^2>E(u)$ (since $E(u)\ge c r^2$ implies $r\le\sqrt{E(u)/c}$). Hence $u(\Sigma)$ is contained in a~ball of radius \smash{$\sqrt{E(u)/c} = \sqrt{C E(u)}$} about $u(\partial\Sigma)$.

 If $E(u)/\bigl(c r_0^2\bigr)>1$, take the maximum number $n$ of disjoint balls of radius $r_0$ that the image of $u$ can meet. Then $n\le E(u)/\bigl(c r_0^2\bigr)$ by monotonicity, and the image is contained in a~ball of radius $(2n+1)r_0\le r_0\bigl(1+2E(u)/\bigl(c r_0^2\bigr)\bigr)<3E(u)/(cr_0) = 3CE(u)/r_0$ around $u(\partial\Sigma)$.
 \end{proof}

For later use, we also record a~localized version, originally observed in~\cite{GromanFloerOpen}.
\begin{df}\label{df:a-bounded}
 For $a>0$, we say that the complete Riemannian metric $g_J$ is \emph{$a$-bounded at a~point} $p\in M$ if \smash{$\inj_{g_J}(p)\ge \frac{1}{a}$} and $\Sec_{g_J}\le a^{2}$ on $B_{1/a}(p)$. We say that $g_J$ is \emph{$a$-bounded on a~set} if it is $a$-bounded at every point of that set.
\end{df}
\begin{rem}
 When $g_J$ is $a$-bounded, we can take $r_0=1/(2a)$ \cite[Theorem 4.4.1 and Remark~4.4.3]{McDuffSalamonJHol}.
\end{rem}

\begin{lm}[local confinement]\label{lem:local-confinement}
Fix $a>0$. There exist $r_0=r_0(a)>0$ and a~universal constant $C>0$ with the following property. Let $K\subset M$ be compact and let $R>0$.
Assume~$J$ is $\omega$-compatible and $g_J$ is complete and is $a$-bounded on the $g_J$-ball $B_{R}(K)$.
Let $u$ be a~$J$-holomorphic curve $u$ with $u(\partial\Sigma)\subset K$ and write
\[
\rho:=\max\bigl(2CE(u)/r_0,\sqrt{CE(u)}\bigr).
\]
If $\rho\le R$, then $u(\Sigma)\subset B_R(K)$.
\end{lm}
\begin{proof}
The proof of Lemma~\ref{lem:sikorav-improved} relies on the monotonicity inequality which depends only on the local geometry of a~complete Riemannian metric. So taking $r_0$ and $C$ as in Lemma~\ref{lem:sikorav-improved}, we obtain by the same argument that
$
u(\Sigma)\subset B_{\rho}(u(\partial\Sigma))
$.
\end{proof}

\begin{rem}
Lemmas~\ref{lem:sikorav}, \ref{lem:sikorav-improved}, and~\ref{lem:local-confinement} also hold for $\omega$-tame almost complex structures $J$ that are geometrically bounded (respectively $a$-bounded on the relevant set) and satisfy the uniform tameness condition $|\omega|_{g_J}\le C$. The constants then depend on the tameness bound as well. See Remark~\ref{rem:omega-tame}.
\end{rem}

\subsection{Proof of Theorem~\ref{tm:1}}\label{subsec:Proof:tm:1}
We now turn to the proof of Theorem~\ref{tm:1}.
The idea is to pick a~small Lagrangian torus $L$ linking the hypersurface~$N$, consider the class of the obvious ``meridian'' holomorphic disc intersecting~$N$ once, and show that the corresponding moduli space has nonzero evaluation degree.

Fix $p\in N$. Choose a~Darboux chart $\psi\colon U\to \bC^{n}$ centered at $p$ such that $\psi(p)=0$ and, in these coordinates,
$
\psi(U\cap N)=\bigl(\{0\}\times \bC^{n-1}\bigr)\cap B_{r}(0)
$
for some $r>0$.
For sufficiently small $t>0$, define a~Lagrangian torus $L_t\subset U$ by
\[
L_t:=\{(z_1,\dots,z_n)\in \bC^n \mid |z_1|=t, |z_i|=10t \text{ for }i=2,\dots,n\}.
\]
Let $\alpha\in \pi_2(M,L_t)$ be the relative homotopy class represented in the chart by the disc
$
u_0\colon  (D,\partial D)\allowbreak\to (\bC^n,L_t)$, $ u_0(\zeta)=(t\zeta,0,\dots,0)$,
so $E(u_0)=\int_D u_0^*\omega_{\mathrm{std}}=\pi t^{2}$.

For an almost complex structure $J$ on $M$, denote by $\mathcal{M}_{J}(L_t,\alpha)$ the moduli space of $J$\nobreakdash-holo\-morphic discs $u\colon(D,\partial D)\to(M,L_t)$ in class $\alpha$ with one boundary marked point, modulo reparametrization. Let
$
\ev\colon \mathcal{M}_{J}(L_t,\alpha)\to L_t
$
be the evaluation map at the marked boundary point.



\begin{lm}\label{lem:degree-one}
There exist arbitrarily small $t > 0$ such that the following holds.
For a~generic geometrically bounded $\omega$-tame almost complex structure $J$ on $M$, the moduli space $\mathcal{M}_{J}(L_t,\alpha)$ is a~smooth closed manifold of dimension $n$, and the evaluation map $\ev\colon\mathcal{M}_{J}(L_t,\alpha)\to L_t$ has degree $1$.
\end{lm}
\begin{proof}
Step 1: \emph{Index and regularity.}
Trivialize $TM$ over the chart and use the standard computation for the product torus to see that the Maslov index of $\alpha$ equals $2$.
Since the disc is simple (in particular not multiply covered), for a~generic choice of $\omega$-tame $J$ the linearized $\bar\partial_J$-operator is surjective (see, e.g., \cite{McDuffSalamonJHol}), hence the regular part $\mathcal{M}_{J,reg}(L_t,\alpha)$ consisting of maps from the disc to $M$ is a~smooth manifold of the expected dimension $n+\mu(\alpha)-2=n$.

Step 2: \emph{Compactness $($no bubbling$)$.}
Since $\omega$ is a~rational form and $H_2(M; \mathbb{Z})$ is finitely generated, the image of $\omega$ on $\pi_2(M)$ is a~discrete subgroup of $\mathbb{R}$, say $\kappa\mathbb{Z}$ for some $\kappa > 0$. Thus there exists $\hbar > 0$ such that every nonconstant $J$-holomorphic sphere has area at least~$\hbar$. Choose a~sufficiently large integer $m$ and define $t > 0$ by the relation $\pi t^2 = \kappa / m$. By choosing~$m$ large enough, we ensure both $\pi t^2 < \hbar$ and $t < r/10$, so the construction stays well inside the Darboux chart. The image of $\omega$ on the relative homotopy group $\pi_2(M, L_t)$ is generated over $\mathbb{Z}$ by the absolute periods on $\pi_2(M)$ and the areas of the basic discs spanning the factors of the torus~$L_t$ in the chart. These generators are $\kappa$, $\pi t^2 = \kappa / m$, and $100\pi t^2 = 100\kappa / m$. Consequently, the area of any relative class in $\pi_2(M, L_t)$ is an integer multiple of $\kappa / m$.
 Any stable map in class $\alpha$ has total area $\pi t^2 = \kappa / m$. It cannot contain a~nonconstant sphere component because~${\kappa / m < \hbar}$. Moreover, any disc bubbling would split $\alpha$ as a~sum of two or more nontrivial relative classes, each requiring a~strictly positive area. Since the minimal positive area available in $\pi_2(M, L_t)$ is exactly~${\kappa / m}$, and the total area of $\alpha$ is $\kappa / m$, such a~split is impossible. Thus~${\mathcal{M}_{J,\text{reg}}(L_t, \alpha) = \mathcal{M}_J(L_t, \alpha)}$ is a~closed manifold.

Step 3: \emph{Independence of $\deg(\ev)$ and normalization.}
Let $J_0$, $J_1$ be geometrically bounded and $\omega$-tame.
By \cite[Theorem 4.15\,(a)]{GromanFloerOpen} (or, alternatively, Theorem~\ref{tm:normalized-completeness}\,(2),~(3)), one can choose a~generic path $\{J_s\}_{s\in[0,1]}$ from $J_0$ to $J_1$ so the moduli space of pairs $\{(s,u)\mid u\in \mathcal{M}_{J_s}(L_t,\alpha)\}$ is contained in a~compact set. Since disc bubbling is excluded as above, this gives a~compact cobordism showing that $\deg(\ev)$ is independent of $J$.

Finally, take $J$ which agrees with the standard complex structure on the Darboux chart and is geometrically bounded on the complement of a~compact set (hence geometrically bounded globally). For $t$ sufficiently small, Lemma~\ref{lem:sikorav} forces any $J$-holomorphic disc in class $\alpha$ with boundary on $L_t$ to remain inside the chart, where the computation reduces to the standard toric case. In that case, regularity and $\deg(\ev)=1$ are well known; see~\cite{ChoOhToric,FOOO_Toric1}.
\end{proof}

We can now prove Theorem~\ref{tm:1}.
\begin{proof}[Proof of Theorem~\ref{tm:1}]
Suppose by contradiction that $M\setminus N$ admits a~geometrically bounded $\omega$-tame almost complex structure $J$, and choose $a>0$ such that $g_J$ is $a$-bounded.
Fix the small parameter $t$ so that Lemma~\ref{lem:degree-one} applies, and note that $\alpha$ has algebraic intersection number $1$ with $N$ (it is the meridian disc in the local model).

Set $E:=E(u_0)=\pi t^2$ and \smash{$R>\max\bigl(2CE/r_0(a),\sqrt{CE}\bigr)$}, with $r_0(a)$ and $C$ from Lemma~\ref{lem:local-confinement} (so that the confinement radius \smash{$\rho=\max\bigl(2CE/r_0,\sqrt{CE}\bigr)\le R$}).
Let $B:=B_{R}(L_t)\subset M\setminus N$ be the closed $g_J$-neighborhood of $L_t$ of radius $R$.
Since $J$ is geometrically bounded on $M\setminus N$, the metric $g_J$ is complete; hence closed bounded subsets are compact, so $B$ is compact.

Choose an $\omega$-tame almost complex structure $J'$ on $M$ which agrees with $J$ on an open neighborhood of $B$ and is geometrically bounded on $M$ (e.g., interpolate to a~fixed geometrically bounded structure outside a~compact set).
By Lemma~\ref{lem:degree-one}, there exists a~stable $J'$-holomorphic disc $u$ in class $\alpha$ with boundary on $L_t$.
Its energy equals $E$.

Since $J'=J$ near $B$ and $g_J$ is $a$-bounded on $B$, Lemma~\ref{lem:local-confinement} implies $u(D)\subset B$. (We use local confinement here because $J'$ is only guaranteed to be $a$-bounded on $B$, not globally on $M$.)
But~${\alpha\cdot N=1}$ forces any representative of $\alpha$ to intersect $N$, contradicting $B\subset M\setminus N$.
\end{proof}

I used AI-assisted tools for editing and improving the exposition (in particular for wording, structure, and catching typographical issues).
All mathematical content, results, and any remaining errors are my own responsibility.

\begin{thebibliography}{99}
\footnotesize\itemsep=0pt

\bibitem{AurouxAnticanonicalTDuality}
Auroux D., Mirror symmetry and {$T$}-duality in the complement of an
 anticanonical divisor, \textit{J.~G\"okova Geom. Topol. GGT} \textbf{1}
 (2007), 51--91, \href{http://arxiv.org/abs/0706.3207}{arXiv:0706.3207}.

\bibitem{ChoOhToric}
Cho C.-H., Oh Y.-G., Floer cohomology and disc instantons of {L}agrangian torus
 fibers in {F}ano toric manifolds,
 \href{https://doi.org/10.4310/AJM.2006.v10.n4.a10}{\textit{Asian~J. Math.}}
 \textbf{10} (2006), 773--814,
 \href{http://arxiv.org/abs/math.SG/0308225}{arXiv:math.SG/0308225}.

\bibitem{EntovGanorMembrez}
Entov M., Ganor Y., Membrez C., Lagrangian isotopies and symplectic function
 theory, \href{https://doi.org/10.4171/CMH/451}{\textit{Comment. Math. Helv.}}
 \textbf{93} (2018), 829--882,
 \href{http://arxiv.org/abs/1610.09631}{arXiv:1610.09631}.

\bibitem{FOOO_Toric1}
Fukaya K., Oh Y.-G., Ohta H., Ono K., Lagrangian {F}loer theory on compact toric
 manifolds.~{I}, \href{https://doi.org/10.1215/00127094-2009-062}{\textit{Duke
 Math.~J.}} \textbf{151} (2010), 23--174,
 \href{http://arxiv.org/abs/0802.1703}{arXiv:0802.1703}.

\bibitem{Gallot}
Gallot S., Hulin D., Lafontaine J., Riemannian geometry, 3rd ed., \textit{Universitext},
 \href{https://doi.org/10.1007/978-3-642-18855-8}{Springer}, Berlin, 2004.

\bibitem{GromanFloerOpen}
Groman Y., Floer theory and reduced cohomology on open manifolds,
 \href{https://doi.org/10.2140/gt.2023.27.1273}{\textit{Geom. Topol.}}
 \textbf{27} (2023), 1273--1390,
 \href{http://arxiv.org/abs/1510.04265}{arXiv:1510.04265}.

\bibitem{GromanWrappedSemiToricSYZ}
Groman Y., The wrapped {F}ukaya category for semi-toric~{C}alabi--{Y}au,
 \href{https://doi.org/10.4171/jems/1466}{\textit{J.~Eur. Math. Soc. (JEMS)}}
 \textbf{26} (2024), 2373--2439,
 \href{http://arxiv.org/abs/1805.03635}{arXiv:1805.03635}.

\bibitem{GromanVarolgunesCompleteEmbeddings}
Groman Y., Varolgunes U., Locality of relative symplectic cohomology for
 complete embeddings,
 \href{https://doi.org/10.1112/s0010437x23007492}{\textit{Compos. Math.}}
 \textbf{159} (2023), 2551--2637,
 \href{http://arxiv.org/abs/2110.08891}{arXiv:2110.08891}.

\bibitem{Gunther}
G{\"u}nther P., Einige S{\"a}tze {\"u}ber das {V}olumen von geod{\"a}tischen
 {K}ugeln, \textit{Ber. Verh. S{\"a}chs. Akad. Wiss. Leipzig Math.-Nat. Kl.}
 \textbf{104} (1960), 1--38.

\bibitem{KarshonTangExcision}
Karshon Y., Tang X., Symplectic excision,
 \href{http://arxiv.org/abs/2101.03534}{arXiv:2101.03534}.

\bibitem{McDuffSymplecticStructures}
McDuff D., Symplectic structures~-- a new approach to geometry, \textit{Notices
 Amer. Math. Soc.} \textbf{45} (1998), 952--960.

\bibitem{McDuffSalamonJHol}
McDuff D., Salamon D., {$J$}-holomorphic curves and symplectic topology, 2nd ed.,
 \textit{Amer. Math. Soc. Colloq. Publ.}, Vol.~52, American
 Mathematical Society, Providence, RI, 2012.

\bibitem{STVFluxShape}
Shelukhin E., Tonkonog D., Vianna R., Geometry of symplectic flux
 and~{L}agrangian torus fibrations,
 \href{https://doi.org/10.1112/topo.70002}{\textit{J.~Topol.}} \textbf{17}
 (2024), e70002, 56~pages,
 \href{http://arxiv.org/abs/1804.02044}{arXiv:1804.02044}.

\bibitem{Sikorav}
Sikorav J.-C., Some properties of holomorphic curves in almost complex
 manifolds, in Holomorphic {C}urves in {S}ymplectic {G}eometry, \textit{Progr.
 Math.}, Vol.~117,
 \href{https://doi.org/10.1007/978-3-0348-8508-9_6}{Birkh\"auser}, Basel,
 1994, 165--189.

\bibitem{StratmannParamRays}
Stratmann B., Removing parametrized rays symplectically,
 \href{https://doi.org/10.4310/jsg.2022.v20.n2.a4}{\textit{J.~Symplectic
 Geom.}} \textbf{20} (2022), 499--508,
 \href{http://arxiv.org/abs/2009.05465}{arXiv:2009.05465}.

\end{thebibliography}
\end{document}
